\documentclass[sigconf]{acmart}

\settopmatter{printacmref=false}
\setcopyright{none}
\renewcommand\footnotetextcopyrightpermission[1]{}

\usepackage{amsmath}
\usepackage{booktabs}
\usepackage{float}

% --------------------------------------------------
% Paper information
% --------------------------------------------------

\title{NLP Assignment 2 Report}

\author{Muhammad Umer Siddiqui}
\email{ms08702@st.habib.pk}
\affiliation{%
  \institution{Dhanani School of Science and Engineering, Habib University}
  \city{Karachi}
  \country{Pakistan}
}

\author{Abdullah Shaikh}
\email{as09245@st.habib.pk}
\affiliation{%
  \institution{Dhanani School of Science and Engineering, Habib University}
  \city{Karachi}
  \country{Pakistan}
}

\author{Arqish Zaria}
\email{az09714@st.habib.pk}
\affiliation{%
  \institution{Dhanani School of Science and Engineering, Habib University}
  \city{Karachi}
  \country{Pakistan}
}

\begin{document}

\maketitle

\section{Classifier Performance}
The Naive Bayes classifier was evaluated on a random test subset of 200 reviews (100 positive, 100 negative). The model achieved an overall accuracy of 82.0\% and an \textbf{F1-score of 0.816}. 

The confusion matrix (Table~\ref{tab:confusion}) demonstrates that the classifier performs consistently across both classes. With a Precision of 0.833 and Recall of 0.800, the model is slightly more conservative in predicting the positive class, ensuring high confidence when it does, though missing a small fraction of actual positive reviews. Overall, the error distribution is well-balanced (20 false negatives vs. 16 false positives), indicating the model did not develop a heavy bias toward either classification despite being trained on a limited 1,000-review subset.

\begin{table}[H]
\caption{Confusion Matrix}
\label{tab:confusion}
\begin{tabular}{@{}lcc@{}}
\toprule
 & \textbf{Predicted Positive} & \textbf{Predicted Negative} \\ \midrule
\textbf{Actual Positive} & 80 & 20 \\
\textbf{Actual Negative} & 16 & 84 \\ \bottomrule
\end{tabular}
\end{table}

\section{Implementation Challenges}
During the implementation and testing of the classifier, several standard Natural Language Processing challenges were encountered and addressed within the code:

\begin{itemize}
    \item \textbf{Floating-Point Underflow:} Multiplying many small probabilities for large text documents causes numerical underflow, yielding a final probability of zero. This was successfully mitigated by operating in the log space, utilizing \texttt{math.log()} and summing the log-probabilities of tokens instead of multiplying raw probabilities.
    \item \textbf{Zero-Frequency Problem:} Words present in the test set that were unseen during training would result in a zero probability, completely nullifying the entire document's score. This was handled using Laplace smoothing (adding 1 to word counts) when calculating the positive and negative likelihoods.
\end{itemize}

\section{Ideas for Improvement}
While the baseline unigram Bag-of-Words model performs well, the following modifications could significantly improve its predictive accuracy:

\begin{itemize}
    \item \textbf{Implementing N-grams:} Transitioning from unigrams to bigrams or trigrams would allow the model to capture local context and crucial semantic modifiers, specifically negations (e.g., distinguishing ``good'' from ``not good'').
    \item \textbf{Advanced Preprocessing:} Incorporating lemmatization or stemming would group morphological variants of words together (e.g., ``running'', ``runs'', and ``ran'' to the root ``run''), thereby condensing the vocabulary space and reinforcing feature weights.
    \item \textbf{TF-IDF Weighting:} Replacing raw frequency counts with Term Frequency-Inverse Document Frequency (TF-IDF) would penalize highly common terms that survived basic stop-word removal, allowing the model to focus strictly on discriminative, class-defining vocabulary.
    \item \textbf{Full Dataset Utilization:} The current model was trained on a randomly sampled subset of 1,000 reviews. Training on the full \texttt{aclImdb} dataset (25,000 training reviews) would dramatically improve the classifier's vocabulary coverage and generalizability.
\end{itemize}

\end{document}